\chapter{Record Bodies and the Xfield Cursor}
\label{ssec:xfield-spec-vs-impl}

An FS-tree leaf record has a fixed-size header and a variable-length
tail of extended fields. The fixed header is easy. The variable tail
is subtle: Apple's reference admits an ambiguous reading, the
converged code is unanimous in a single direction, and the right rule
must be applied uniformly across every record family that uses
xfields.

\section{The Cursor Rule}

\textbf{Observation.} The xfield tail follows one cursor rule, derived
from \texttt{linux-apfs-rw}, \texttt{apfs-fuse}, \texttt{dissect.apfs},
\texttt{go-apfs}, \texttt{libfsapfs}, and the Sleuth Kit:

\begin{enumerate}
  \item The tail begins with \texttt{xf\_blob\_t}: 4 bytes containing
    \texttt{xf\_num\_exts} and \texttt{xf\_used\_data}.
  \item Immediately after is the metadata table:
    \texttt{xf\_num\_exts} entries of \texttt{x\_field\_t}, each 4
    bytes (\texttt{x\_type}, \texttt{x\_flags}, \texttt{x\_size}).
  \item The value area begins immediately after the metadata table.
    There is no extra alignment of the value-area start.
  \item Reading in metadata-table order, each value consumes
    \texttt{round\_up(x\_size, 8)} bytes. The cursor advances by the
    padded length, not by \texttt{x\_size}.
  \item \texttt{xf\_used\_data} equals
    \texttt{sum(round\_up(x\_size, 8))} for the record's xfields.
\end{enumerate}

That is the rule. It is one rule across all xfield-bearing record
families, and a record whose \texttt{xf\_used\_data} disagrees with
the padded sum is malformed.

\begin{figure}[h]
\centering
\begin{tikzpicture}[
  every node/.style={font=\footnotesize},
  byte/.style={draw, minimum width=0.7cm, minimum height=0.75cm,
               align=center, inner sep=1pt},
  blob/.style={byte, fill=gray!15},
  meta/.style={byte, fill=blue!15},
  val/.style={byte, fill=green!18},
  pad/.style={byte, fill=red!15, font=\scriptsize\itshape},
  brace/.style={decorate, decoration={brace, amplitude=4pt,
                                       raise=1pt, mirror}},
]
  % blob header
  \node[blob] (b0) {};
  \node[blob, right=0pt of b0] (b1) {};
  \node[blob, right=0pt of b1] (b2) {};
  \node[blob, right=0pt of b2] (b3) {};

  % metadata table -- 3 entries
  \node[meta, right=0.06cm of b3] (m0) {};
  \foreach \i [evaluate=\i as \prev using {int(\i-1)}] in {1,...,11} {
    \node[meta, right=0pt of m\prev] (m\i) {};
  }

  % value 1: x_size=11, padded to 16
  \node[val, right=0.06cm of m11] (v0) {};
  \foreach \i [evaluate=\i as \prev using {int(\i-1)}] in {1,...,10} {
    \node[val, right=0pt of v\prev] (v\i) {};
  }
  \foreach \i [evaluate=\i as \prev using {int(\i-1)}] in {11,...,15} {
    \node[pad, right=0pt of v\prev] (v\i) {p};
  }

  % value 2: x_size=40 (we show a compressed view)
  \node[val, right=0.06cm of v15, minimum width=2.1cm] (val2) {value 2 (40\,B)};

  % value 3: x_size=8 (padded to 8 -- no pad needed)
  \node[val, right=0.06cm of val2, minimum width=0.95cm] (val3) {8};

  % labels
  \node[above=0.05cm of b1.east, anchor=south, font=\scriptsize]
    {\texttt{xf\_blob\_t}};
  \node[above=0.05cm of m5.east, anchor=south, font=\scriptsize]
    {metadata: 3 $\times$ 4 B};
  \node[above=0.05cm of v3.east, anchor=south, font=\scriptsize]
    {value 1 (11\,B + 5\,B padding)};

  % cursor marker
  \draw[-Stealth, thick] ($(m11.south east) + (0, -0.55)$)
    -- ($(m11.south east) + (0, -0.05)$)
    node[below=0.55cm, font=\scriptsize, align=center]
      {cursor starts here\\(no extra alignment)};
\end{tikzpicture}
\caption{The xfield tail for the sparse inode at OID 23
  (\texttt{xf\_num\_exts=3}, sizes 11/40/8). The value cursor starts
  immediately after the metadata table; each value advances by
  \texttt{round\_up(x\_size, 8)} bytes. Total padded value-area length
  is $16+40+8 = 64$, matching \texttt{xf\_used\_data}.}
\end{figure}

The rule fits in a few lines of code:

\begin{minted}{rust}
let mut cursor = metadata_end;
let mut padded_total = 0usize;
for entry in metadata.iter() {
    let value = &blob[cursor .. cursor + entry.x_size as usize];
    let padded = (entry.x_size as usize + 7) & !7;
    cursor += padded;
    padded_total += padded;
}
assert_eq!(padded_total, blob_header.xf_used_data as usize);
\end{minted}

\section{Why The Reference Reads Otherwise}

Apple's reference describes the xfield layout in terms that admit a
"the value area is independently eight-byte aligned" reading. Under
that reading, the cursor would jump to the next eight-byte boundary
after the metadata table before reading the first value. The
difference is small and frequently invisible: when
\texttt{xf\_num\_exts * 4} is already a multiple of eight, both
readings agree.

The converged code is uniform that no such re-alignment happens. The
value cursor starts at the metadata-table end, full stop. Each
value's padding to the next eight-byte multiple is performed
\emph{after} the value, not before the first.

The project initially modelled four candidate layouts in EX-13,
scoring each against the proof fixture and selecting per-record.
That probe successfully recovered the correct values but did not
explain why. SR-015 reviewed the converged sources and showed that
the four candidates were a probe artefact: a single rule explains
every record the candidates fit. EX-16 re-ran the decode under the
single rule and confirmed it.

\section{Evidence From The Proof Fixture}

\textbf{Observation.} On the proof fixture, 14 records carry
xfield tails. Under the single cursor rule, every one of them
satisfies \texttt{xf\_used\_data == sum(round\_up(x\_size, 8))}, and
every one of them reconstructs the namespace and per-file logical
size identically to the mounted POSIX oracle.

The records and their xfield shapes:

\begin{center}
\begin{tabular}{rlrlrr}
\toprule
oid & family & num\_exts & x\_size set & xf\_used\_data & padded sum \\
\midrule
2  & inode   & 1 & 5                 & 8  & 8 \\
3  & inode   & 1 & 12                & 16 & 16 \\
16 & inode   & 1 & 4                 & 8  & 8 \\
17 & inode   & 1 & 4                 & 8  & 8 \\
17 & dir\_rec & 1 & 8                & 8  & 8 \\
17 & dir\_rec & 1 & 8                & 8  & 8 \\
18 & inode   & 1 & 11                & 16 & 16 \\
19 & inode   & 2 & 15, 40            & 56 & 56 \\
20 & inode   & 2 & 10, 40            & 56 & 56 \\
23 & inode   & 3 & 11, 40, 8         & 64 & 64 \\
24 & inode   & 2 & 40, 10            & 56 & 56 \\
26 & inode   & 1 & 9                 & 16 & 16 \\
27 & inode   & 2 & 17, 40            & 64 & 64 \\
28 & inode   & 2 & 17, 40            & 64 & 64 \\
\bottomrule
\end{tabular}
\end{center}

The sparse inode at OID 23 is the load-bearing row. It carries three
xfields (\texttt{NAME=11}, \texttt{DSTREAM=40},
\texttt{SPARSE\_BYTES=8}) with padded lengths summing to 64; the
\texttt{NAME} value has a non-multiple-of-eight \texttt{x\_size} (11),
which forces the cursor to use padded advancement.  Earlier
candidate-scoring had needed per-record reasoning for this row;
under the single rule it decodes cleanly.

No record has unused trailing bytes inside its xfield blob. The
structural assertion holds without exception.

\section{The Required Xfield Types}

The xfield universe is large; v1 needs four types and tolerates
others.

\begin{description}
  \item[\texttt{INO\_EXT\_TYPE\_NAME}.] The inode name. The name is
    also reachable through DIR\_REC keys, and v1 prefers DIR\_REC
    names for namespace construction (they carry the directory
    placement). The inode-name xfield is a consistency check.
  \item[\texttt{INO\_EXT\_TYPE\_DSTREAM}.] A \texttt{j\_dstream\_t}
    embedded in the inode, carrying logical size, allocated size,
    default crypto id, and a default-extent reference.
  \item[\texttt{INO\_EXT\_TYPE\_SPARSE\_BYTES}.] A count of bytes the
    sparse-allocation layer has not committed to disk. This is not a
    logical size; chapter~8 develops the distinction.
  \item[\texttt{DREC\_EXT\_TYPE\_SIBLING\_ID}.] A directory record's
    sibling-group identifier for hard links. SIBLING\_MAP closes
    the mapping from sibling ID to inode.
\end{description}

Other types (\texttt{INO\_EXT\_TYPE\_FS\_UUID},
\texttt{INO\_EXT\_TYPE\_PURGEABLE\_FLAGS}, and so on) appear in the
wild but are not consumed by v1 product rows. The cursor rule still
reads past them correctly because the rule treats every value as
opaque bytes of length \texttt{x\_size} padded to a multiple of
eight; only the values the parser actually inspects need a known
type.

\section{Field-Level Cross-Tool Equality}

The xfield rule, the fixed-header decoders, and the fail-closed gates
of the next chapter together specify the body decoder completely.
The project verified the specification by running two independent
decoders against the same raw container and comparing every field.

\textbf{Observation.} On the proof fixture, the Rust body decoder
emits 53 records. An independent Python decoder, built from the
same SR-015 rule and the same fail-closed boundary, emits 53
records. The two record lists key identically on
\texttt{(node\_paddr, entry\_index)}. Per-record field comparison
finds zero divergent fields across \texttt{key}, \texttt{value}, every
\texttt{xfields[*]}, \texttt{xfield\_used\_data},
\texttt{xfield\_padded\_total},
\texttt{xfield\_unused\_trailing\_bytes}, \texttt{validation\_notes},
\texttt{key\_len}, \texttt{value\_len}, \texttt{object\_id},
\texttt{raw\_type}, and \texttt{family}.

Cross-tool field equality is the strongest oracle the body decoder
will see until the project has a same-run \texttt{go-apfs} body dump
to add. The Python decoder and the Rust decoder are independent
implementations of the same source-review rules; their agreement
makes the rules themselves testable.

\section{When To Use This Chapter}

A reader implementing an APFS body decoder should:

\begin{enumerate}
  \item read the fixed header for the record type from Apple's
    reference,
  \item apply the cursor rule above to the xfield tail,
  \item validate against every fail-closed condition in chapter~7,
  \item compare against an independent reference implementation, not
    against intuition.
\end{enumerate}

Steps 3 and 4 are not optional. The next chapter is what they look
like in practice.
